% MAT670 -- Interlude F (placed last, after Interlude E): computation and certified proof.
% Entirely new material (not in the handwritten notes). Labels prefixed COMP:.
% Scripts referred to in the text live in /home/claude/mat670/code/.

\chapter[Interlude F: Computation and proof]{Interlude F: Computation and certified proof}
\label{COMP:ch}

\begin{addedblock}
This interlude is new: it is not in the handwritten notes. It collects the computational ideas that
sit behind several results quoted in the course (the Kepler conjecture, the Tammes problem for $13$
and $14$ points, the kissing numbers in dimensions $4$, $8$ and $24$) and tries them out on small
problems that the notes already contain. To keep it readable it is set in black; all of it should be
read as added material. The short Python programs mentioned below are in the directory
\texttt{code/} next to these notes.
\end{addedblock}

Look again at the proof of Thue's theorem in Lecture~1. It begins with an infinite object, a
saturated packing of the plane, and it ends with a statement about a single triangle
(Lemma~\ref{L1:lem:triangle-density}, in the form its proof actually uses): a triangle with sides at
least $2$ and circumradius less than $2$ has area at least $\sqrt3$. Everything in between is bookkeeping: a Delaunay triangulation cuts the
packing into pieces, a weighted average reduces the density of the whole to the density of the pieces,
and the boundary terms disappear in the limit. The Fejes T\'oth inequality of Lecture~2 has the same
shape. A triangulation of the sphere is cut into $2n-4$ spherical triangles, Euler's formula counts
them, and one inequality about spherical triangles finishes the job.

This is the shape of almost every sharp result in the subject.
\begin{enumerate}[label=(\arabic*)]
\item \emph{Reduce} a continuum of configurations to finitely many combinatorial \emph{cases}, indexed by
  a geometric object: a triangulation, a contact graph, a Delaunay star, a perfect form.
\item \emph{Prune} the cases by cheap necessary conditions: angle sums, areas, Euler's formula, planarity.
\item \emph{Settle} each surviving case by an inequality (it cannot occur, or it is worse than the
  candidate) or by solving equations (it is the candidate).
\item \emph{Certify} every step, so that the proof can be checked by somebody who does not trust the
  person, or the program, that found it.
\end{enumerate}
In Lecture~1 there was one case and the inequality took two lines. In Hales's proof of the Kepler
conjecture there are thousands of cases and tens of thousands of inequalities. This interlude is about
what happens in between: how cases are indexed and enumerated (Section~\ref{COMP:sec:index}), how
searches are organized (Section~\ref{COMP:sec:bb}), which pieces of linear algebra and optimization do
the work (Sections~\ref{COMP:sec:matrix}--\ref{COMP:sec:kkt}), what algebra can and cannot decide
(Section~\ref{COMP:sec:poly}), and how a computer can produce a \emph{proof} rather than a number
(Section~\ref{COMP:sec:verify}). In Section~\ref{COMP:sec:thue} we put the pieces together and prove
Thue's triangle lemma by a certified computer search.

%======================================================================
\section{The shape of a computer-assisted proof}
\label{COMP:sec:shape}

Here are the problems from the course, arranged by the four steps above.
\begin{center}
\renewcommand{\arraystretch}{1.25}
\small
\begin{tabular}{p{0.19\textwidth}p{0.23\textwidth}p{0.27\textwidth}p{0.18\textwidth}}
\toprule
problem & cases indexed by & each case settled by & reference\\
\midrule
disks in $\R^2$ & Delaunay triangles & area $\ge\sqrt3$ & Lecture~1\\
$N$ caps on $\Sph^2$ & triangulations, Euler & spherical area inequality & Lecture~2\\
critical points of $\rho$ & balanced contact graphs & force balance (an LP) & Lectures~10--11\\
Tammes, $N{=}13,14$ & irreducible graphs & LP and inequalities & Musin--Tarasov\\
kissing, $n=8,24$ & --- (one case) & one polynomial & Section~\ref{COMP:sec:matrix}\\
Kepler, $\R^3$ & tame plane graphs & nonlinear inequalities, LPs & Hales\\
lattices in $\R^n$ & perfect forms & Voronoi's algorithm & Lect.~6, Int.~B\\
\bottomrule
\end{tabular}
\end{center}
Two remarks before we start. First, the line between ``by hand'' and ``by computer'' is not where one
might think. The kissing number $240$ in dimension $8$ is proved by a single polynomial, which a
computer found but which anyone can check with pencil and paper (Example~\ref{COMP:ex:240}). On the
other hand, the enumeration of all tame graphs for the Kepler conjecture cannot be checked by hand at
all. Second, what matters in the end is not whether a computer was used but whether every step
produced a \emph{certificate}: a finite object (a list of graphs, a polynomial with rational
coefficients, a collection of boxes, a vector of weights) together with a simple procedure that checks
it. Most of this interlude is about which certificates are available and how to produce them.

%======================================================================
\section{Indexing cases by geometric objects}
\label{COMP:sec:index}

\subsection{Contact graphs}
Lectures~10 and~11 attached to a configuration $X=(x_1,\dots,x_N)$ on $\Sph^2$ with minimal distance
$\theta=\rho(X)$ its \emph{contact graph} (Definition~\ref{L11:def:contact}), with edges the geodesic arcs of
length $\theta$. Two simple lemmas turn this graph into a good index for case analysis.

\begin{lemma}[Contact graphs are planar]\label{COMP:lem:planar}
Let $0<\theta<\pi$ and let $X\subset\Sph^2$ have all pairwise distances $\ge\theta$. Then two contact
edges (arcs of length $\theta$ between points of $X$) can meet only at a common endpoint.
\end{lemma}

\begin{proof}
Two minor arcs from a common endpoint $a$ lie on great circles through $a$, which meet again only at
$-a$, so they meet only at $a$. Suppose the arcs $[a,b]$ and $[c,d]$ have four distinct endpoints and
cross at a point $p$. By the triangle inequality,
\[
  d(a,c)+d(b,d)\le d(a,p)+d(p,c)+d(b,p)+d(p,d)=d(a,b)+d(c,d)=2\theta .
\]
Both terms on the left are $\ge\theta$, so both equal $\theta$ and both triangle inequalities are
equalities. Since $d(a,c)=\theta<\pi$, equality in $d(a,c)\le d(a,p)+d(p,c)$ means that $p$ lies on the
minor arc from $a$ to $c$. That arc and $[a,b]$ leave $a$ in the same direction (towards $p$) and both
have length $\theta$, so $c=b$, a contradiction.
\end{proof}

\begin{lemma}[Angle lemma]\label{COMP:lem:angle}
Let $u,v,w\in\Sph^2$ with $d(u,v)=d(v,w)=\theta\in(0,\pi)$ and $d(u,w)\ge\theta$. Then the angle
$\varphi$ at $v$ between the arcs $vu$ and $vw$ satisfies $\varphi\ge\alpha(\theta)$, where
\[
  \cos\alpha(\theta)=\frac{\cos\theta}{1+\cos\theta}
\]
is the angle of the equilateral spherical triangle of side $\theta$. Consequently a point of $X$ with
$k\ge2$ contacts has $k\,\alpha(\theta)\le2\pi$.
\end{lemma}

\begin{proof}
By the spherical law of cosines, $\cos d(u,w)=\cos^2\theta+\sin^2\theta\cos\varphi$. The hypothesis
$\cos d(u,w)\le\cos\theta$ gives
$\cos\varphi\le(\cos\theta-\cos^2\theta)/\sin^2\theta=\cos\theta/(1+\cos\theta)$. The $k$ angles between
consecutive contact arcs at a point add up to $2\pi$.
\end{proof}

For $\theta=\pi/3$ this gives $\cos\alpha=\frac13$: \emph{at most five contacts}, because
$5\arccos\frac13<2\pi<6\arccos\frac13$. This is the same number $\arccos\frac13\approx70.53^\circ$ as the
dihedral angle of the regular tetrahedron in Section~\ref{TET:sec}, and for the same reason. For
$\theta=\pi/2$ we get $\alpha=\pi/2$ and at most four contacts.

\subsection{A worked example: the critical configurations of five points}
Exercise~\ref{L10:ex:classify} of Lecture~10 asked for the classification of balanced configurations,
and the notes did $N=3$ and $N=4$ by hand. Let us do $N=5$ the way a computer would.

Let $X$ be a critical configuration of $\rho$ on $\Conf(5,\Sph^2)$, with balancing weights $w\ge0$, not
all zero, as in Theorem~\ref{L11:thm:forcebalance}. Let $S$ be the \emph{stressed graph}: the edges with
$w_e>0$, together with all five points (points without stressed edges are isolated vertices of $S$;
they feel no force). We look for all graphs that can occur as $S$, using four rules.
\begin{enumerate}[label=\textbf{R\arabic*},leftmargin=*]
\item $S$ is planar. (Lemma~\ref{COMP:lem:planar}, since $\theta\le\max\rho<\pi$.)
\item No vertex has degree $1$. (A single positive force cannot balance.)
\item No vertex of degree $2$ lies in a triangle of $S$. (If $v$ has exactly two stressed edges, the two
  forces are opposite, so the neighbours $u,w$ lie on a great circle on either side of $v$ and
  $d(u,w)=\min(2\theta,2\pi-2\theta)$. If also $d(u,w)=\theta$, then $\theta=2\pi/3$. But four unit vectors
  with pairwise inner products $\le-\frac12$ would have $\abs{x_1+\dots+x_4}^2\le4-12\cdot\frac12<0$.)
\item If two vertices $a,b$ have three common neighbours, then $a,b$ are not adjacent, and their
  common neighbours span no triangle. (The circles of radius $\theta$ about $a$ and $b$ meet in at most two
  points unless they coincide, which forces $b=-a$ and $\theta=\pi/2$. Then $d(a,b)=\pi\ne\theta$, and the
  common neighbours lie on the equator between $a$ and $b$, where three points cannot be pairwise
  $\pi/2$ apart.)
\end{enumerate}
The program \texttt{code/contact5.py} runs through the $34$ graphs on five vertices (from the
\emph{Atlas of Graphs} built into the \texttt{networkx} library). One is empty; R1 removes $K_5$; R2 removes $18$;
R3 removes $7$ and R4 removes one more. Six graphs survive (Figure~\ref{COMP:fig:N5}).

\begin{figure}[ht]
\centering
\begin{tikzpicture}[scale=0.95,
  v/.style={circle,fill,inner sep=1.15pt}, f/.style={circle,draw,fill=white,inner sep=1.15pt},
  lab/.style={font=\scriptsize,align=center}]
  % row 1
  \begin{scope}[shift={(0,0)}]
    \node[v] (a) at (0,0.6) {}; \node[v] (b) at (-0.55,-0.35) {}; \node[v] (c) at (0.55,-0.35) {};
    \draw (a)--(b)--(c)--(a); \node[f] at (-0.7,0.55) {}; \node[f] at (0.7,0.55) {};
    \node[lab] at (0,-0.95) {$C_3$\\ R3};
  \end{scope}
  \begin{scope}[shift={(2.1,0)}]
    \node[v] (a) at (-0.45,0.45) {}; \node[v] (b) at (0.45,0.45) {}; \node[v] (c) at (0.45,-0.45) {}; \node[v] (d) at (-0.45,-0.45) {};
    \draw[addcol,thick] (a)--(b)--(c)--(d)--(a); \node[f] at (0,0) {};
    \node[lab] at (0,-0.95) {$C_4$\\ \textbf{realized}, $\pi/2$};
  \end{scope}
  \begin{scope}[shift={(4.2,0)}]
    \node[v] (a) at (0,0.6) {}; \node[v] (b) at (-0.45,0) {}; \node[v] (c) at (0.45,0) {}; \node[v] (d) at (0,-0.6) {};
    \draw (a)--(b)--(d)--(c)--(a) (b)--(c); \node[f] at (0.75,0.55) {};
    \node[lab] at (0,-0.95) {diamond\\ R3};
  \end{scope}
  \begin{scope}[shift={(6.3,0)}]
    \foreach \k in {0,...,4} \node[v] (p\k) at (90+72*\k:0.6) {};
    \draw[addcol,thick] (p0)--(p1)--(p2)--(p3)--(p4)--(p0);
    \node[lab] at (0,-0.95) {$C_5$\\ \textbf{realized}, $2\pi/5$};
  \end{scope}
  \begin{scope}[shift={(8.4,0)}]
    \node[v] (a) at (0,0.62) {}; \node[v] (b) at (-0.58,-0.38) {}; \node[v] (c) at (0.58,-0.38) {}; \node[v] (d) at (0,0) {};
    \draw (a)--(b)--(c)--(a) (d)--(a) (d)--(b) (d)--(c); \node[f] at (0.72,0.55) {};
    \node[lab] at (0,-0.95) {$K_4$\\ free point};
  \end{scope}
  \begin{scope}[shift={(10.5,0)}]
    \node[v] (o) at (0,0) {}; \node[v] (a) at (-0.6,0.42) {}; \node[v] (b) at (-0.6,-0.42) {};
    \node[v] (c) at (0.6,0.42) {}; \node[v] (d) at (0.6,-0.42) {};
    \draw (o)--(a)--(b)--(o)--(c)--(d)--(o);
    \node[lab] at (0,-0.95) {bowtie\\ R3};
  \end{scope}
  \begin{scope}[shift={(12.6,0)}]
    \node[v] (a) at (-0.4,-0.5) {}; \node[v] (b) at (0.4,-0.5) {}; \node[v] (c) at (0.4,0.12) {}; \node[v] (d) at (-0.4,0.12) {};
    \node[v] (e) at (0,0.62) {};
    \draw (a)--(b)--(c)--(d)--(a) (c)--(e)--(d);
    \node[lab] at (0,-0.95) {house\\ R3};
  \end{scope}
  % row 2
  \begin{scope}[shift={(0,-2.5)}]
    \node[v] (n) at (0,0.6) {}; \node[v] (s) at (0,-0.6) {};
    \node[v] (a) at (-0.6,0) {}; \node[v] (b) at (0,0) {}; \node[v] (c) at (0.6,0) {};
    \draw[addcol,thick] (n)--(a)--(s) (n)--(b)--(s) (n)--(c)--(s);
    \node[lab] at (0,-0.95) {$K_{2,3}$\\ \textbf{realized}, $\pi/2$};
  \end{scope}
  \begin{scope}[shift={(2.1,-2.5)}]
    \node[v] (n) at (0,0.6) {}; \node[v] (s) at (0,-0.6) {};
    \node[v] (a) at (-0.6,0) {}; \node[v] (b) at (0.28,0) {}; \node[v] (c) at (0.72,0) {};
    \draw (n)--(a)--(s) (n)--(b)--(s) (n)--(c)--(s) (n)--(s);
    \node[lab] at (0,-0.95) {$K_{1,1,3}$\\ R3};
  \end{scope}
  \begin{scope}[shift={(4.2,-2.5)}]
    \node[v] (t) at (0,0.6) {}; \node[v] (a) at (-0.72,-0.3) {}; \node[v] (b) at (-0.25,-0.45) {};
    \node[v] (c) at (0.25,-0.45) {}; \node[v] (d) at (0.72,-0.3) {};
    \draw (a)--(b)--(c)--(d) (t)--(a) (t)--(b) (t)--(c) (t)--(d);
    \node[lab] at (0,-0.95) {gem\\ R3};
  \end{scope}
  \begin{scope}[shift={(6.3,-2.5)}]
    \node[v] (n) at (0,0.6) {}; \node[v] (s) at (0,-0.6) {};
    \node[v] (a) at (-0.6,0) {}; \node[v] (b) at (0,0) {}; \node[v] (c) at (0.6,0) {};
    \draw (n)--(a)--(s) (n)--(b)--(s) (n)--(c)--(s) (a)--(b);
    \node[lab] at (0,-0.95) {$K_{2,3}+e$\\ weights};
  \end{scope}
  \begin{scope}[shift={(8.4,-2.5)}]
    \node[v] (a) at (0,0.62) {}; \node[v] (b) at (-0.55,-0.3) {}; \node[v] (c) at (0.55,-0.3) {}; \node[v] (d) at (0,0.02) {};
    \node[v] (e) at (0,-0.68) {};
    \draw (a)--(b)--(c)--(a) (d)--(a) (d)--(b) (d)--(c) (e)--(b) (e)--(c);
    \node[lab] at (0,-0.95) {$K_4$ + ear\\ R3};
  \end{scope}
  \begin{scope}[shift={(10.5,-2.5)}]
    \node[v] (a) at (-0.45,0.45) {}; \node[v] (b) at (0.45,0.45) {}; \node[v] (c) at (0.45,-0.45) {}; \node[v] (d) at (-0.45,-0.45) {};
    \node[v] (o) at (0,0) {};
    \draw (a)--(b)--(c)--(d)--(a) (o)--(a) (o)--(b) (o)--(c) (o)--(d);
    \node[lab] at (0,-0.95) {$W_4$\\ weights};
  \end{scope}
  \begin{scope}[shift={(12.6,-2.5)}]
    \node[v] (a) at (-0.62,-0.4) {}; \node[v] (b) at (0.62,-0.4) {}; \node[v] (c) at (0,0.62) {};
    \node[v] (o) at (0,0) {}; \node[v] (e) at (0,-0.75) {};
    \draw (a)--(b)--(c)--(a) (o)--(a) (o)--(b) (o)--(c) (e)--(a) (e)--(b);
    \draw (e) .. controls (1.15,-0.8) and (0.95,0.6) .. (c);
    \node[lab] at (0,-0.95) {$K_5-e$\\ R4};
  \end{scope}
\end{tikzpicture}
\caption{The $14$ graphs on five vertices that pass R1 and R2 (hollow: a point without stressed
contacts). R3 and R4 remove eight of them. Of the six that remain, three are removed by a one-line
geometric argument and three (blue) are realized by critical configurations.}
\label{COMP:fig:N5}
\end{figure}

Now the last, geometric, step, which here is easy enough to do by hand.
\begin{itemize}
\item $K_4$ plus a point: four pairwise equidistant points form a regular tetrahedron, so
  $x_1+x_2+x_3+x_4=0$. For any fifth point $y$ the four numbers $\ip y{x_i}$ add up to $0$, so one of them is
  $\ge0$ and $d(y,x_i)\le\pi/2<\theta$.
\item $W_4$: by R4 applied to the two pairs of opposite rim vertices, $\theta=\pi/2$, the rim is a square on a
  great circle and the hub is its pole. At a rim vertex the two rim forces are opposite and the force
  towards the hub is perpendicular to both, so its weight must vanish.
\item $K_{2,3}+e$: the same argument at an endpoint of the extra edge (north and south forces opposite, the
  extra force along the equator).
\item The survivors are realized: $C_5$ by the regular pentagon on a great circle ($\theta=2\pi/5$); $C_4$
  by a square on a great circle with the fifth point at a pole ($\theta=\pi/2$); $K_{2,3}$ by the two
  poles and three points on the equator with all gaps in $[\pi/2,\pi)$, as in Figure~\ref{TOP:fig:N5} of
  Interlude~D. For $C_5$ and $C_4$ equal weights balance. For $K_{2,3}$ put the same weight $w_k$ on both
  edges at the $k$-th equatorial point (so the equatorial points balance); since all gaps are less than
  $\pi$, the three directions at a pole are not in a closed half-plane, so there are $w_1,w_2,w_3>0$
  balancing the poles. (Equal weights work only for gaps $2\pi/3$.)
\end{itemize}

\begin{proposition}\label{COMP:prop:N5}
The critical values of $\rho$ on $\Conf(5,\Sph^2)$ are $2\pi/5$ and $\pi/2$. At $2\pi/5$ the critical
configurations are the regular pentagons on great circles; at $\pi/2$ they are the configurations with
two antipodal points and three points on the equator between them, with all gaps $\ge\pi/2$.
\end{proposition}

Together with Lecture~11 this is the complete list for $N\le5$. The lowest critical value is $2\pi/N$,
as predicted by the Kusner--Kusner--Lagarias--Shlosman theorem quoted there. For $N=6$ the same rules
leave many more graphs, and the geometric step needs actual computation: one has to solve, for each
graph, the equations ``all edges have length $\theta$, all non-edges are longer, the forces balance''
(Section~\ref{COMP:sec:poly}). That is the general pattern.

\subsection{The same idea at full scale}
\paragraph{Tammes for $N=13,14$.} The Tammes problem asks for $\max\rho$. Musin and Tarasov call a
contact graph \emph{irreducible} (or jammed) if no single point can be shifted so as to increase its
distance to all its neighbours and no ``Danzer flip'' is possible, and they show that for $N>6$ the
contact graph of an optimal configuration is irreducible. The
contact graph of an irreducible configuration is planar, every vertex has degree $0$, $3$, $4$ or $5$
(the upper bound is Lemma~\ref{COMP:lem:angle}), every face is a convex spherical polygon with sides $\theta$ and at most
$\lfloor2\pi/\theta\rfloor$ vertices, and isolated vertices lie in faces with at least six vertices.
These are the analogues of R1--R4. They enumerate the planar graphs with these properties by computer,
discard most of them by area and angle inequalities (in the spirit of the proof of the Fejes T\'oth
inequality: the face areas add up to $4\pi$, and each face has a minimal area depending on its
shape), and treat the few survivors by optimization. This gives the optimal $13$- and $14$-point
configurations (Section~\ref{COMP:sec:gallery}). Generating planar graphs with prescribed degrees and face
sizes is itself a well-developed subject; the standard tool is \texttt{plantri} (Brinkmann and McKay).

\paragraph{Kepler.} In Hales's proof the cases are indexed by \emph{tame plane graphs}. Around each ball
one draws, on a sphere about its centre, the planar graph whose vertices are (the radial projections of)
the neighbouring centres within a fixed distance ($2.51$ in the original proof, $2.52$ in the revised and formal
one) and whose edges join neighbours that are close to each other. The local density (the \emph{score} of a
\emph{decomposition star}, later replaced by \emph{Marchal cells}) is then bounded case by case.
``Tame'' is a list of combinatorial conditions of the R1--R4 type (degree bounds, face sizes, a weight
inequality), designed so that a non-tame graph is excluded by an easy argument. Hales's original archive
had $5128$ tame graphs. When Nipkow, Bauer and Schultz re-did the enumeration inside the proof assistant
Isabelle (2006), they confirmed that the archive was complete and found it redundant, reducing it to
$2771$ graphs. The revised proof uses a larger archive of $18\,762$ graphs.

\paragraph{Lattices.} For lattice packings the index is a \emph{perfect form}
(Theorem~\ref{NT:gal:voronoi} of Interlude~B). Voronoi's algorithm walks from one perfect form to its
neighbours across the facets of a polyhedral cone, so the enumeration is a graph search. Its cost is
dominated by symmetry: one must decide whether a new form is equivalent to one already found. In
dimension $8$ there are $10\,916$ perfect forms (Dutour Sikiri\'c, Sch\"urmann and Vallentin, 2007).

%======================================================================
\section{Branch and bound, and meet in the middle}
\label{COMP:sec:bb}

\subsection{Branch and bound}
Suppose we want to show that a function $F$ on a compact set $K$ satisfies $F\ge0$, or to find
$\max_KF$. \emph{Branch and bound} keeps a list of pieces of $K$ (boxes, cells, partial assignments).
For each piece $B$ it computes a cheap \emph{bound}: a number $\underline F(B)\le\inf_BF$ (for the
inequality) or $\overline F(B)\ge\sup_BF$ (for the maximum). If the bound settles the piece, the
piece is discarded; otherwise it is \emph{branched}, that is, split into smaller pieces. Two
properties make it work: the bounds must be \emph{valid} (never wrong), and they must become
\emph{sharp} as the pieces shrink.

\paragraph{Max-min problems.} The min-type function $\rho$ of Lecture~10 is well suited to this. Let
$B_1,\dots,B_N\subset\Sph^2$ be small cells and suppose $x_i\in B_i$. Then
\[
  \rho(x_1,\dots,x_N)\ \le\ \overline\rho(B_1,\dots,B_N):=\min_{i<j}\ \max_{y\in B_i,\,z\in B_j}d(y,z),
\]
and the right side is a minimum of $\binom N2$ numbers that are easy to compute. If
$\overline\rho<\rho(X_0)$ for a known configuration $X_0$, no configuration with $x_i\in B_i$ can be
optimal. The difficulty is the number of pieces: with cells of size $\varepsilon$ there are about
$(4\pi/\varepsilon^2)^N$ choices, and symmetry multiplies the work by $\abs{\SO(3)\times\Sigma_N}$-many
copies of the same configuration. For the analogous problem of $n$ equal circles in a square,
Mark\'ot and Csendes made this approach work with interval arithmetic and a clever elimination of
regions (the ``method of active areas''), proving the optimal packings for $n=28,29,30$
(about $50$, $50$ and $20$ hours of computing in 2005); Mark\'ot later did $n=31,32,33$ (2021). For
the Tammes problem with $13$ points, the graph enumeration of Section~\ref{COMP:sec:index} is used
instead: fixing the combinatorics first removes both the symmetry and most of the dimension.

\paragraph{Branch and bound on integer points.} The same idea finds short vectors in a lattice. Let $G$ be
a positive definite Gram matrix and $C>0$; we want all $x\in\Z^n$ with $x^TGx\le C$. Completing squares
(the $LDL^T$ decomposition, in exact rational arithmetic) writes
\[
  x^TGx=\sum_{i=1}^nq_{ii}\Big(x_i+\sum_{j>i}q_{ij}x_j\Big)^2,\qquad q_{ii}>0 .
\]
Choose $x_n$ first: the $i=n$ term alone must be $\le C$, which leaves finitely many integers. Given
$x_n$, the remaining budget bounds $x_{n-1}$, and so on. This is the algorithm of Fincke and Pohst
(1985), a branch and bound whose pieces are the partial vectors $(x_k,\dots,x_n)$. Run on the Cartan
matrix of $E_8$ (\texttt{code/fincke\_pohst.py}), it visits $739$ nodes and finds the $240$ roots; with
$C=4$ it visits $5477$ nodes and finds the $2400$ nonzero vectors of norm at most $4$. The size of the
tree depends heavily on the basis: a basis with long, nearly parallel vectors gives a very unbalanced
tree. The LLL algorithm of Lenstra, Lenstra and Lov\'asz (1982) computes, in polynomial time, a reduced
basis whose first vector is within a factor $2^{(n-1)/2}$ of the shortest, and in practice LLL followed by
Fincke--Pohst enumeration is how minimal vectors are found. Voronoi's algorithm needs exactly this at
every step: the minimal vectors of each new form.

\subsection{Meet in the middle}
The idea of \emph{meet in the middle} is to split a search over pairs $(u,v)$ into two half-searches and
match the halves through a table. The classic example is \emph{subset sum}: given $n$ integers, is there
a subset with sum $T$? Split the integers into two halves, list the $2^{n/2}$ subset sums of each half, sort
one list, and look up $T-s$ for every $s$ in the other. This takes about $2^{n/2}$ steps instead of $2^n$
(Horowitz and Sahni, 1974), and Schroeppel and Shamir (1981) reduced the memory to about $2^{n/4}$.

Here is the same idea counting vectors of $E_8$, as defined in Lecture~6
(Definition~\ref{L04b:def:E8}). Write $x=(u,v)$ with $u,v\in\R^4$. Then $x\in E_8$ exactly when $u,v$ are
both integral or both in $(\Z+\frac12)^4$, and the coordinate sum is even. So it suffices to tabulate the
half-vectors $u$ by the pair $\big(\abs u^2,\ \text{sum of $u$ modulo }2\big)$, separately for the two
cosets, and to multiply the counts of matching classes. The program \texttt{code/e8\_mitm.py} lists only
$625$ half-vectors and recovers
\[
  240,\quad 2160,\quad 6720,\quad 17520\qquad\text{vectors of norm }2,4,6,8,
\]
which is $240\,\sigma_3(m)$ as predicted by the theta series (equation~\eqref{NT:eq:E8theta} in
Interlude~B), whereas the full list has $26\,641$ vectors. In lattice language: $E_8$ is a union of four cosets of
$D_4\oplus D_4$, and the theta series of an orthogonal sum is the product of the theta series. This
``gluing'' is how Conway and Sloane describe many lattices.

To be honest about where meet in the middle is used: as an \emph{algorithm} it is central in coding
theory and cryptography. Stern's method (1989) for finding low-weight codewords splits an information set
in two halves and looks for collisions, and Howgrave-Graham (2007) combined lattice reduction with a
meet-in-the-middle search to attack the NTRU cryptosystem. In the proofs about packings quoted in this
course it appears mainly in the gluing form above, as a way to count or list vectors without listing the
whole space.

%======================================================================
\section{Matrix techniques}
\label{COMP:sec:matrix}

\subsection{Gram matrices}
A configuration of $N$ unit vectors is determined, up to an orthogonal transformation, by its Gram
matrix $G_{ij}=\ip{x_i}{x_j}$.

\begin{proposition}\label{COMP:prop:gram}
A real symmetric $N\times N$ matrix $G$ is the Gram matrix of $N$ vectors in $\R^d$ if and only if $G$ is
positive semidefinite of rank at most $d$. The vectors are then unique up to $\Orth(d)$.
\end{proposition}

\begin{proof}
If $G=V^TV$ with $V$ a $d\times N$ matrix, then $y^TGy=\abs{Vy}^2\ge0$ and $\operatorname{rank}G\le d$.
Conversely, diagonalize $G=Q\Lambda Q^T$ with $\Lambda\ge0$ having at most $d$ nonzero entries, and take
$V$ to be the nonzero rows of $\Lambda^{1/2}Q^T$. If $V^TV=W^TW$, the map $Vy\mapsto Wy$ is a well defined
isometry between the column spaces, and it extends to an element of $\Orth(d)$.
\end{proof}

So ``is there a configuration of $N$ points on $\Sph^{d-1}$ with all angles at least $\theta$?'' becomes ``is
there a positive semidefinite matrix with ones on the diagonal, off-diagonal entries $\le\cos\theta$, and
rank at most $d$?'' Without the rank condition this is a \emph{semidefinite program}. The rank condition
is what makes the problem hard, and dropping it is what makes bounds computable. Sometimes, however, the
sign pattern alone settles everything.

\begin{proposition}[Rankin 1955]\label{COMP:prop:rankin}
Let $x_1,\dots,x_N\in\R^d$ be nonzero.
\begin{enumerate}[label=(\alph*),nosep]
\item If $\ip{x_i}{x_j}<0$ for all $i\ne j$, then $N\le d+1$.
\item If $\ip{x_i}{x_j}\le0$ for all $i\ne j$, then $N\le2d$.
\end{enumerate}
\end{proposition}

\begin{proof}
\emph{Claim:} vectors with pairwise non-positive inner products that lie in an open half-space
$\{\ip ye>0\}$ are linearly independent. Indeed, a linear relation can be written
$u=\sum_{i\in P}a_ix_i=\sum_{j\in Q}b_jx_j$ with $a,b\ge0$ and $P\cap Q=\emptyset$. Then
$\abs u^2=\sum a_ib_j\ip{x_i}{x_j}\le0$, so $u=0$, and $0=\ip ue=\sum a_i\ip{x_i}e$ forces $a=0$; similarly
$b=0$. For (a), apply the claim to $x_1,\dots,x_{N-1}$ and $e=-x_N$. For (b), choose $e$ with
$\ip{x_i}e\ne0$ for all $i$ and apply the claim to the vectors on each side.
\end{proof}

Since $d(x,y)\ge\pi/2$ on $\Sph^2$ means $\ip xy\le0$, this settles two cases of the Tammes problem at
once and bounds a third: five or six points cannot all be more than $\pi/2$ apart (by (a), $N\le4$), and
the triangular bipyramid and the octahedron show that $\pi/2$ is attained for $N=5,6$; seven points cannot
all be $\pi/2$ apart (by (b), $N\le6$), so the Tammes distance for $N=7$ is less than $\pi/2$ (its actual
value, $\approx77.87^\circ$, needs much more work).

\subsection{Positive definite kernels and the linear programming bound}
The proof of Proposition~\ref{COMP:prop:rankin} used one positive semidefinite matrix. The linear
programming bound uses infinitely many at once. For $S^{n-1}\subset\R^n$ let $G_k$ be the Gegenbauer
polynomial $C_k^{(n-2)/2}$, normalized by $G_k(1)=1$. By the addition theorem for spherical harmonics,
$G_k(\ip xy)$ is a positive multiple of $\sum_mY_m(x)Y_m(y)$, summed over an orthonormal basis of the
harmonic polynomials of degree $k$. Hence for every finite $C\subset S^{n-1}$
\begin{equation}\label{COMP:eq:posdef}
  \sum_{x,y\in C}G_k(\ip xy)\ \ge\ 0,
\end{equation}
that is, the matrix $\big(G_k(\ip xy)\big)_{x,y\in C}$ is positive semidefinite (Schoenberg, 1942).

\begin{theorem}[Delsarte--Goethals--Seidel 1977; Kabatiansky--Levenshtein 1978]\label{COMP:thm:LP}
Let $f(t)=\sum_{k=0}^Dc_kG_k(t)$ with $c_0>0$ and $c_k\ge0$ for all $k$, and suppose $f(t)\le0$ for
$-1\le t\le\cos\theta$. Then every $C\subset S^{n-1}$ with all angles at least $\theta$ has
$\abs C\le f(1)/c_0$.
\end{theorem}

\begin{proof}
Sum $f(\ip xy)$ over all pairs. The off-diagonal terms are $\le0$, so the sum is at most $\abs Cf(1)$. By
\eqref{COMP:eq:posdef} every $k\ge1$ contributes $\ge0$ and $k=0$ contributes $c_0\abs C^2$. Hence
$c_0\abs C^2\le\abs Cf(1)$.
\end{proof}

This is the third time in these notes that a bound comes from \emph{throwing away squares}: Blichfeldt
dropped $\abs{\sum x_i}^2$ (Lemma~\ref{L2:lem:f0}), Cohn and Elkies dropped one square for each
frequency (Remark~\ref{NT:rem:blichfeldt-vs-CE}), and here each $G_k$ hides a sum of squares.

\begin{example}[The kissing number in dimension $8$]\label{COMP:ex:240}
Take $n=8$, $\theta=\pi/3$, and the polynomial of Odlyzko--Sloane and Levenshtein,
\[
  f(t)=(t+1)\Big(t+\frac12\Big)^2t^2\Big(t-\frac12\Big).
\]
Clearly $f\le0$ on $[-1,\frac12]$. Expanding in the $G_k$ for $n=8$ (exact rational arithmetic; a few lines of
\texttt{sympy}) gives
\[
  (c_0,\dots,c_6)=\Big(\frac3{320},\ \frac3{40},\ \frac{15}{64},\ \frac{39}{80},\ \frac{399}{640},\ \frac9{16},\
  \frac{33}{128}\Big),
\]
all positive, and $f(1)/c_0=\frac94\cdot\frac{320}3=240$. So $\kiss_8\le240$, and $E_8$ gives equality.
The same computation with
$f(t)=(t+1)(t+\frac12)^2(t+\frac14)^2t^2(t-\frac14)^2(t-\frac12)$ (degree $10$) in dimension $24$ gives
$196\,560$, with all eleven coefficients $c_0,\dots,c_{10}$ positive. Note the roots: they are the inner products $-1,-\frac12,0,\frac12$
occurring in the $E_8$ kissing configuration, and $-1,\pm\frac12,\pm\frac14,0$ in the Leech one. Equality in
the proof forces $f(\ip xy)=0$ for $x\ne y$, which is how one guesses $f$, and how Bannai and Sloane
(1981) proved that both kissing configurations are unique.
\end{example}

The polynomial was found by solving a linear program (maximize $c_0$ subject to $f(1)=1$ and to
sampled constraints $f(t_\ell)\le0$), but the \emph{certificate} is the polynomial: rational coefficients,
checked exactly. In dimensions $3$ and $4$ the LP bound is not sharp (it gives $13$ and $25$). Musin
proved $\kiss_4=24$ (2008) by a modification of the method. Replacing pairs by triples of points gives
\emph{semidefinite} programming bounds (Bachoc and Vallentin, 2008), which improved the upper bounds in
dimensions $5,6,7,9,10$; Leijenhorst and de Laat (2024) improved them in dimensions $11$ to $23$. For
packings in $\R^n$ the continuous analogue is the Cohn--Elkies bound (Theorem~\ref{NT:thm:CE}), and its
three-point refinement (Cohn, de Laat and Salmon, 2022) improves the density bounds in dimensions $4$--$7$
and $9$--$16$.

\subsection{Rigidity matrices and stresses}
The other matrix in this subject is the \emph{rigidity matrix} $A$ of a contact graph: the linear map that
sends a motion $v=(v_1,\dots,v_N)$ to the rates of change of the contact distances (Lecture~11, and
Interlude~E, Proposition~\ref{SM:prop:jam}). Its kernel is the space of infinitesimal motions, and the
kernel of $A^T$ is the space of \emph{stresses}, systems of edge weights that balance at every vertex.
Theorem~\ref{L11:thm:forcebalance} says that a configuration is critical when $\ker A^T$ meets the
non-negative orthant. Two computational remarks. First, if the configuration has algebraic coordinates,
ranks and kernels can be computed exactly over a number field, and a positive stress, once found
numerically, can be checked exactly. Second, deciding whether $\ker A^T$ contains a non-negative nonzero
vector is a linear program: maximize nothing subject to $A^Tw=0$, $w\ge0$, $\sum w_e=1$. Its feasibility
is exactly Gordan's alternative from Lecture~10. The next section explains why this is no accident.

%======================================================================
\section{KKT: forces are Lagrange multipliers}
\label{COMP:sec:kkt}

The Tammes problem $\max\rho$ is a nonsmooth optimization problem, since $\rho$ is a minimum of smooth
functions. The standard trick is to add a variable. Write $s$ for the cosine of the minimal distance and
consider
\begin{equation}\label{COMP:eq:tammes}
  \text{minimize } s\quad\text{subject to}\quad \ip{x_i}{x_j}\le s\ \ (i<j),\qquad \abs{x_i}^2=1\ \ (1\le i\le N).
\end{equation}
Now everything is smooth (indeed polynomial), and a configuration $X$ is a local maximum of $\rho$ if and
only if $(X,\cos\rho(X))$ is a local minimum of \eqref{COMP:eq:tammes}. The basic necessary condition for
smooth constrained problems is due to F.~John (1948) for inequality constraints, and to Mangasarian and
Fromovitz (1967) with equality constraints.

\begin{theorem}[Fritz John conditions]
Let $p$ be a local minimum of $f$ subject to $g_k\le0$ and $h_l=0$, all $C^1$. Then there are
$\lambda_0\ge0$, $\lambda_k\ge0$ and $\mu_l\in\R$, not all zero, with
\[
  \lambda_0\nabla f(p)+\sum_k\lambda_k\nabla g_k(p)+\sum_l\mu_l\nabla h_l(p)=0,\qquad \lambda_kg_k(p)=0\ \text{for all }k.
\]
If $\lambda_0>0$ one may take $\lambda_0=1$; these are the \emph{Karush--Kuhn--Tucker (KKT)} conditions.
\end{theorem}

Usually one needs a \emph{constraint qualification} (for instance, linear independence of the gradients of
the active constraints) to guarantee $\lambda_0>0$. For the Tammes problem it comes for free, and the
multipliers turn out to be the forces of Lecture~11.

\begin{proposition}\label{COMP:prop:kkt}
A configuration $X$ with $\rho(X)=\theta\in(0,\pi)$ satisfies the KKT conditions for
\eqref{COMP:eq:tammes} (with $s=\cos\theta$) if and only if it is critical in the sense of
Theorem~\ref{L11:thm:forcebalance}. The multipliers of the constraints $\ip{x_i}{x_j}\le s$ are balancing
weights on the contact graph. In particular every local maximum of $\rho$ is critical.
\end{proposition}

\begin{proof}
With $g_{ij}=\ip{x_i}{x_j}-s$ and $h_i=\abs{x_i}^2-1$, the derivative of the Lagrangian in $s$ and in $x_i$
gives
\[
  \lambda_0=\sum_{i<j}\lambda_{ij},\qquad \sum_{j\ne i}\lambda_{ij}x_j+2\mu_ix_i=0\quad(1\le i\le N),
\]
with $\lambda_{ij}=\lambda_{ji}$, and $\lambda_{ij}=0$ unless $\ip{x_i}{x_j}=s$, that is, unless $ij$ is a
contact. For a contact, $x_j=\cos\theta\,x_i+\sin\theta\,u_{ij}$, with $u_{ij}$ the unit tangent vector at
$x_i$ pointing to $x_j$. The tangential part of the second equation is therefore
$\sin\theta\sum_j\lambda_{ij}u_{ij}=0$: the $\lambda_{ij}$ balance the contact graph. If all
$\lambda_{ij}$ vanished, then $\lambda_0=0$ and $2\mu_ix_i=0$, so all multipliers would vanish; hence the
weights are not all zero and $\lambda_0=\sum\lambda_{ij}>0$. Conversely, balancing weights $w_{ij}$, scaled
so that $\sum w_{ij}=1$, give KKT multipliers $\lambda_{ij}=w_{ij}$, $\lambda_0=1$ and
$\mu_i=-\frac12\cos\theta\sum_jw_{ij}$ (the normal part of the second equation). The last sentence follows
from the Fritz John theorem.
\end{proof}

So the dictionary is: \emph{multipliers are forces}, \emph{complementary slackness says that only contacts
carry force}, and \emph{Gordan's alternative is LP duality}. The weights are found by a linear program,
and the forces at each vertex close up into a polygon (Figure~\ref{COMP:fig:forces}); that polygon is the
first step of the Maxwell--Cremona correspondence between stresses and reciprocal diagrams.

\begin{figure}[ht]
\centering
\begin{tikzpicture}[scale=1.5]
  \begin{scope}
    \fill (0,0) circle (1.6pt) node[above left,font=\small] {$x_i$};
    \foreach \a/\w/\l in {80/0.684/1,190/1/2,330/1/3} {
      \draw[gray] (0,0)--(\a:1.35); \fill (\a:1.35) circle (1.3pt) node[font=\scriptsize,shift={(\a:0.22)}] {$x_{\l}$};
      \draw[-{Stealth},flagcol,very thick] (0,0)--(\a:\w);}
    \draw[gray,dashed] (0,0)--(265:1.35); \fill (265:1.35) circle (1.3pt) node[font=\scriptsize,below] {$x_4$};
    \node[font=\scriptsize,align=center] at (0,-1.85) {weights $0.68,\ 1,\ 1$ on the solid contacts;\\ the dashed contact carries no force};
  \end{scope}
  \begin{scope}[xshift=3.6cm,yshift=-0.2cm]
    \coordinate (P0) at (0,0); \coordinate (P1) at (0.119,0.674); \coordinate (P2) at (-0.866,0.5);
    \draw[-{Stealth},flagcol,very thick] (P0)--(P1); \draw[-{Stealth},flagcol,very thick] (P1)--(P2);
    \draw[-{Stealth},flagcol,very thick] (P2)--(P0);
    \fill (P0) circle (1pt);
    \node[font=\scriptsize,align=center] at (-0.3,-1.65) {the three forces, head to tail,\\ close up: $\sum_jw_{ij}u_{ij}=0$};
  \end{scope}
\end{tikzpicture}
\caption{Force balance at one vertex of a critical configuration, seen in the tangent plane at $x_i$. The
multipliers of Proposition~\ref{COMP:prop:kkt} are the lengths of the arrows. A fourth contact with zero
weight is allowed: the weights are $\ge0$, not necessarily $>0$.}
\label{COMP:fig:forces}
\end{figure}

\begin{remark}[Necessary, not sufficient]
KKT points need not be maxima: the square of Lecture~11 is critical but not a local maximum of $\rho$. A
sufficient condition comes from Interlude~E. If $X$ is \emph{first-order jammed} (every nontrivial
infinitesimal motion strictly shortens some contact, which by Proposition~\ref{SM:prop:jam} means a positive
stress on every contact plus infinitesimal rigidity modulo rotations), then a compactness argument on the
unit sphere of motions shows that $X$ is a strict local maximum of $\rho$ up to rotation. The
$\rho$-evolution algorithm of Lecture~11 is, in this language, an \emph{active set method}: it moves in a
direction that increases all active constraints until a new constraint becomes active, and it stops at a
KKT point.
\end{remark}

%======================================================================
\section{Polynomial methods}
\label{COMP:sec:poly}

\subsection{Equations for balanced configurations}
Once a contact graph is fixed, a critical configuration is a solution of polynomial equations:
$\abs{x_i}^2=1$, $\ip{x_i}{x_j}=c$ on the edges, and the balance equations
$\sum_jw_{ij}(x_j-c\,x_i)=0$, modulo rotations, together with the inequalities $\ip{x_i}{x_j}<c$ off the
edges and $w\ge0$. Polynomial algebra solves the equations exactly.

\begin{example}[The square antiprism, $N=8$]
The Tammes configuration for $8$ points (Sch\"utte and van der Waerden, 1951) is a square antiprism with
all $16$ edges equal. Put the two squares at heights $\pm h$ on circles of radius $r$, rotated by
$45^\circ$ against each other. A side of a square has squared length $2r^2$ and a slanted edge has
$(2-\sqrt2)r^2+4h^2$, and $\cos\theta=1-r^2$. With $s=\sqrt2$ and $r_2=r^2$, the equations are
\[
  r_2+h^2-1=0,\qquad s^2-2=0,\qquad s\,r_2-4h^2=0,\qquad c-1+r_2=0 .
\]
A Gr\"obner basis in the lexicographic order $r_2>h>s>c$ (Buchberger's algorithm, one line in
\texttt{sympy}) is
\[
  \{\,r_2+c-1,\ \ h^2-c,\ \ 2s-7c-1,\ \ 7c^2+2c-1\,\}.
\]
It is triangular, like the result of Gaussian elimination, and its last element is a polynomial in $c$
alone. So $c=\cos\theta$ satisfies $7c^2+2c-1=0$, that is $\cos\theta=(2\sqrt2-1)/7$ and
$\theta\approx74.8585^\circ$.
\end{example}

For larger systems one needs to decide how many real roots a univariate polynomial has in an
interval. \emph{Sturm's theorem} (1829) does it exactly: with $p_0=p$, $p_1=p'$ and
$p_{k+1}=-\operatorname{rem}(p_{k-1},p_k)$, the number of roots in $(a,b]$ is the drop in the number of sign
changes of $(p_0,p_1,\dots)$ from $a$ to $b$. For $p=7c^2+2c-1$ the sequence is $(7c^2+2c-1,\,14c+2,\,\frac87)$,
with signs $(-,+,+)$ at $0$ and $(+,+,+)$ at $1$: exactly one root in $(0,1]$. Descartes' rule of signs
and its refinements (Vincent; Collins and Akritas, 1976) give faster \emph{real root isolation}.
\emph{Resultants} eliminate one variable at a time, and \emph{homotopy continuation} (as in the programs
Bertini and PHCpack) follows the solutions of an easy system along a path to all isolated complex solutions
of the given one, numerically. Numerical solutions can be certified afterwards: by Smale's
\emph{$\alpha$-theory}, if $\alpha(f,x)=\beta(f,x)\gamma(f,x)<(13-3\sqrt{17})/4\approx0.1577$, where
$\beta=\norm{Df(x)^{-1}f(x)}$ is the Newton step and
$\gamma=\sup_{k\ge2}\norm{Df(x)^{-1}D^kf(x)/k!}^{1/(k-1)}$, then Newton's method started at $x$ converges
quadratically to a true zero. The program \texttt{alphaCertified} (Hauenstein and Sottile, 2012) checks
these conditions in exact or interval arithmetic.

\subsection{Tarski--Seidenberg: everything is decidable}
A \emph{first-order formula over $\R$} is built from polynomial equations and inequalities with rational
coefficients, using $\wedge,\vee,\neg$ and quantifiers $\exists x,\forall x$ over real numbers. For example,
``there are $13$ points on $\Sph^2$ with pairwise distances at least $\pi/3$'' is
\[
  \exists x_1,\dots,x_{13}\in\R^3:\ \bigwedge_i\abs{x_i}^2=1\ \wedge\ \bigwedge_{i<j}\ip{x_i}{x_j}\le\tfrac12 .
\]

\begin{theorem}[Tarski 1948/1951; Seidenberg 1954]
Every first-order formula over $\R$ is equivalent to a quantifier-free one, and the equivalent formula can
be computed. In particular the truth of every first-order sentence over $\R$ is decidable.
\end{theorem}

So $\kiss_3\le12$ could, \emph{in principle}, be decided by running an algorithm. There is also a
nice theoretical consequence. Apply quantifier elimination to the formula in one free variable
$c$: ``there are $N$ unit vectors with $\ip{x_i}{x_j}\le c$''. The result defines a finite union of
points and intervals with real algebraic endpoints, and its least element is the cosine of the Tammes
distance. \emph{Hence every Tammes value $\cos d_N$ is an algebraic number}, as are the optimal
values of all the finite configuration problems of the course, without anyone having to find them.

Why not just do it? Collins's \emph{cylindrical algebraic decomposition} (1975), the first practical
algorithm, has running time roughly $(sd)^{2^{O(k)}}$ for $s$ polynomials of degree $d$ in $k$ variables,
\emph{doubly} exponential in $k$, and Davenport and Heintz (1988) showed that quantifier elimination
can be inherently doubly exponential. For a single block of existential quantifiers there are better
algorithms (Grigoriev--Vorobjov, Renegar, Basu--Pollack--Roy) with running time of the order
$s^{k+1}d^{O(k)}$, singly exponential. For the $13$-sphere sentence above, $k=39$ and $s=91$; that is far
beyond any computer. In practice, cylindrical algebraic decomposition handles problems with a handful of
variables. Tarski's theorem decides everything in principle and nothing big in practice. The art of
the subject is to reduce a big problem to many small ones, each with few variables, and then to treat
those by the methods of the next section.

%======================================================================
\section{Verification}
\label{COMP:sec:verify}

\subsection{Interval arithmetic}
\emph{Interval arithmetic} (R.~E.~Moore, 1966) computes with closed intervals $X=[\underline x,\overline x]$
instead of numbers:
\[
  X+Y=[\underline x+\underline y,\ \overline x+\overline y],\qquad
  X\cdot Y=[\min P,\max P],\quad P=\{\underline x\underline y,\underline x\overline y,\overline x\underline y,\overline x\overline y\},
\]
and so on. On a computer each lower endpoint is rounded down and each upper endpoint up (the IEEE~754
standard provides these \emph{directed rounding} modes), or one uses exact rational endpoints, as we will.
By induction on the length of an expression one gets the \emph{fundamental theorem}: if an expression $F$ is
evaluated on intervals, the result contains $F(x)$ for every choice of $x$ in the intervals. So a single
evaluation on a box gives a \emph{proved} enclosure of the range of $F$ on the box.

The price is overestimation, the \emph{dependency problem}: each occurrence of a variable is treated as
independent. For $f(x)=x-x^2$ on $X=[0,1]$, the natural evaluation $X-X\cdot X=[0,1]-[0,1]=[-1,1]$, while the
true range is $[0,\frac14]$. Rewriting $f=\frac14-(x-\frac12)^2$ gives the exact range. The \emph{centred
form} $f(m)+f'(X)(X-m)$, here $\frac14+[-1,1]\cdot[-\frac12,\frac12]=[-\frac14,\frac34]$, is still not exact,
but its overestimation is $O(w^2)$ in the width $w$ of the box, against $O(w)$ for the natural evaluation.
This is why bounds become sharp as boxes shrink, which is what branch and bound needs.

\paragraph{Certified zeros.} Interval arithmetic also proves that solutions exist.

\begin{proposition}[Interval Newton]\label{COMP:prop:newton}
Let $f\in C^1[a,b]$, $X=[a,b]$, $m\in X$, and suppose that the interval $f'(X)$ (any enclosure of $f'$ on
$X$) does not contain $0$. Put $N(X)=m-f(m)/f'(X)$. If $N(X)\subseteq X$, then $f$ has exactly one zero in $X$,
and it lies in $N(X)$.
\end{proposition}

\begin{proof}
Since $f'\ne0$ on $X$, $f$ is strictly monotone, so there is at most one zero. Say $f'>0$ and $f(m)\ge0$ (the
other cases are symmetric). By the mean value theorem $f(a)=f(m)-f'(\xi)(m-a)$ for some $\xi\in X$, and
$m-f(m)/f'(\xi)\in N(X)\subseteq X$ gives $f(m)/f'(\xi)\le m-a$, so $f(a)\le0$. By the intermediate value
theorem there is a zero $x^*$ in $[a,m]$. Finally $x^*=m-f(m)/f'(\eta)$ for some $\eta$ between $x^*$ and $m$,
so $x^*\in N(X)$.
\end{proof}

For the antiprism, take $f(c)=7c^2+2c-1$, $X=[0.26,0.27]$ and $m=0.265$. Then $f(m)=0.021575$,
$f'(X)=[5.64,5.78]$, and $N(X)=[0.26117\ldots,0.26126\ldots]\subset X$, so the Tammes cosine for $8$ points
lies in $N(X)$. One more step shrinks the interval to width $4\cdot10^{-9}$: Newton's quadratic convergence,
now with a proof attached. In several variables the same role is played by the \emph{Krawczyk operator}
$K(X)=m-Yf(m)+(I-Yf'(X))(X-m)$, with $Y$ an approximate inverse of $f'(m)$: if $K(X)$ lies in the interior of
$X$, then $f$ has a unique zero in $X$ (Krawczyk 1969; Moore 1977). This is how a numerically computed
balanced configuration is turned into a theorem.

\subsection{Dual numbers and jets}
To use centred forms or interval Newton one needs enclosures of derivatives on boxes. They come almost for
free from \emph{dual numbers}: compute in $\R[\varepsilon]/(\varepsilon^2)$. For a polynomial,
\[
  f(a+b\varepsilon)=f(a)+f'(a)\,b\,\varepsilon,
\]
because $\varepsilon^2=0$ kills everything else. So evaluating $f$ at $x+\varepsilon$ with the ordinary rules
$(a+a'\varepsilon)(b+b'\varepsilon)=ab+(ab'+a'b)\varepsilon$ produces $f(x)$ and $f'(x)$ together. This is
\emph{forward-mode automatic differentiation}. With $k$ variables one carries a gradient
$(a,a_1',\dots,a_k')$ instead; with $\R[\varepsilon]/(\varepsilon^{r+1})$ one carries the Taylor coefficients
up to order $r$, a \emph{jet}. Now let the entries be intervals. Evaluating on a box $B$ then gives, in one
pass, enclosures of $f(B)$ and of $\nabla f(B)$, by the fundamental theorem. We use exactly this in
Section~\ref{COMP:sec:thue}.

Going one step further, a \emph{Taylor model} (Berz and Makino, 1998) represents a function on a box by a
polynomial of degree $r$ together with an interval bounding the remainder. Arithmetic on Taylor models
keeps the dependency in the polynomial part, so the overestimation is $O(w^{r+1})$. The nonlinear inequalities
of the Kepler conjecture were proved with second-order versions of this idea (``Taylor intervals''), and
Solovyev and Hales (2013) implemented them inside the proof assistant HOL Light.

\subsection{Sums of squares}
A polynomial $p$ that is a sum of squares of polynomials is visibly nonnegative, and the squares are a
certificate. Write $p=m^TQm$ with $m$ the vector of monomials of degree at most $\frac12\deg p$. Then $p$ is
a sum of squares if and only if this holds for some positive semidefinite $Q$, since $Q=L^TL$ gives
$p=\abs{Lm}^2$. The conditions on $Q$ are linear equations (matching coefficients) plus $Q\succeq0$:
\emph{finding a sum-of-squares certificate is a semidefinite program} (Parrilo, 2000/2003; Lasserre, 2001).

Not every nonnegative polynomial is a sum of squares (Hilbert, 1888). Motzkin's
$x^4y^2+x^2y^4-3x^2y^2+1$ is nonnegative by the inequality between arithmetic and geometric means applied
to $x^4y^2,x^2y^4,1$, but it is not a sum of squares. To prove $p\ge0$ only on a set
$K=\{g_1\ge0,\dots,g_r\ge0\}$ one looks for certificates of the form
\[
  p=\sigma_0+\sum_i\sigma_ig_i\qquad(\text{Putinar, 1993}),\qquad\text{or}\qquad
  p=\sum_{I\subseteq\{1,\dots,r\}}\sigma_I\prod_{i\in I}g_i\qquad(\text{Schm\"udgen, 1991}),
\]
with sums of squares $\sigma$. Schm\"udgen showed that a polynomial strictly positive on a compact $K$ always
has the second kind of certificate, and Putinar that, under an extra hypothesis on $K$ (Archimedean), it has the
first. The Positivstellensatz of Krivine (1964) and Stengle (1974) gives certificates for the
\emph{emptiness} of a semialgebraic set, and so for nonnegativity in general, at the price of denominators.
Compare: the $\sigma_i$ are polynomial versions of the KKT multipliers of Section~\ref{COMP:sec:kkt}, and the
multiplier $\lambda(z)$ of Section~\ref{COMP:sec:thue} is an example.

The semidefinite program is solved numerically, so the matrix $Q$ it returns is only approximately positive
semidefinite and only approximately satisfies the linear equations. Peyrl and Parrilo (2008) showed how to
round it to an \emph{exact} rational certificate when $p$ is strictly inside the cone: project the rounded $Q$
onto the affine space of exact solutions and check $Q\succeq0$ exactly, for example by an exact Cholesky
factorization. The same rounding idea turns the numerical solutions of the LP and SDP bounds of
Section~\ref{COMP:sec:matrix} into rigorous bounds.

\subsection{Formal proof}
A certificate still has to be checked by \emph{some} program, and programs have bugs. A \emph{proof
assistant} (HOL Light, Isabelle, Coq, Lean) reduces this risk to a small trusted kernel that checks every
logical step.
\begin{itemize}
\item \emph{Kepler.} Hales's proof was published in 2005 (with a joint paper of Hales and Ferguson). The
  refereeing could not certify all of the computer code, and Hales started the Flyspeck project to produce a
  formal proof. It was completed on 10 August 2014, in HOL Light and
  Isabelle. The formal proof checks more than $23\,000$ nonlinear inequalities (after subdividing domains),
  which took about $5000$ processor hours on the Microsoft Azure cloud, together with $43\,078$ linear
  programs and the tame graph enumeration of Nipkow, Bauer and Schultz.
\item \emph{Dimension $8$.} Viazovska's proof (Interlude~B) relies on two inequalities between
  modular forms. In her paper they were checked with computer help: the functions were replaced by truncated
  expansions at the two ends of the interval, the truncation errors were bounded rigorously, and the
  inequalities for the truncations were checked by a finite calculation. Romik (2023) later gave a proof
  without computer calculations. A Lean formalization of the whole dimension-$8$ theorem, begun by Hariharan
  and Viazovska in 2024, was completed in February 2026; its final stages were produced by the
  autoformalization system Gauss of Math, Inc.\ (Hariharan, Birkbeck, Lee, Ma, Mehta, Poiroux and
  Viazovska, arXiv:2604.23468). According to EPFL's announcement (March 2026) the dimension-$24$ case followed.
\item \emph{Outside this course.} The four colour theorem (Appel and Haken, 1976--77; a simpler proof by
  Robertson, Sanders, Seymour and Thomas, 1997) was formalized in Coq by Gonthier (2005, published 2008).
  Tucker's proof (1999, 2002) that the Lorenz equations have a strange attractor, Smale's 14th problem,
  is a classic of interval analysis: it combines rigorous integration of the flow with a normal form near
  the origin.
\end{itemize}

%======================================================================
\section{A certified proof of Thue's triangle lemma}
\label{COMP:sec:thue}

We now prove the inequality of Lemma~\ref{L1:lem:triangle-density} by computer, using branch and bound, exact interval
arithmetic, dual numbers and a KKT multiplier. Of course Lecture~1 proves it in two lines. The point is to
watch every piece of the machinery work on an example where we know the answer.

\paragraph{The polynomial form.} Let a triangle have sides $a,b,c\ge2$ and circumradius $R\le2$, and let
$c$ be the longest side. Put $x=a^2$, $y=b^2$, $z=c^2$. Heron's formula says
\[
  16\operatorname{Area}^2=H(x,y,z):=2xy+2yz+2zx-x^2-y^2-z^2,
\]
and $R=abc/(4\operatorname{Area})$ gives $R^2=xyz/H$, so $R\le2$ means $g:=4H-xyz\ge0$. Since $c\le2R\le4$,
the domain is
\[
  D=\{4\le x\le z,\ 4\le y\le z,\ z\le16,\ g\ge0\},
\]
and the inequality of Lemma~\ref{L1:lem:triangle-density} (area $\ge\sqrt3$, i.e.\ density
$\le\pi/\sqrt{12}$) follows from $H\ge48$ on $D$. (We allow $R=2$; this only makes the claim stronger.) All transcendental functions are gone: this is
the first rule of thumb.

\paragraph{Where it is tight.} $H=48$ at the equilateral triangle $(4,4,4)$, a corner of $D$. It is also
$48$ at $(4,4,12)$, the triangle with sides $2,2,2\sqrt3$ and largest angle $2\pi/3$, which has $R=2$
exactly. (It is excluded by the strict inequality $R<2$ in Lecture~1, but it lies in the closure, so the
inequality on $D$ is sharp at two points.) Tight points are the enemy of branch and bound: near them no bound
of the form ``$H\ge48$ on this box'' can be proved with any slack, and a naive search never ends.

\paragraph{The KKT multiplier.} At $(4,4,12)$ the constraint $g\ge0$ is active. Solving the KKT
equations $\nabla H=\mu_x e_x+\mu_ye_y+\lambda\nabla g$ there gives $\nabla H=(24,24,-8)$,
$\nabla g=(48,48,-48)$, hence $\lambda=\frac16$ and $\mu_x=\mu_y=16$. At $(4,4,4)$ the constraint $g$ is not
active ($g=128$) and its multiplier is $0$. On the edge $x=y=4$ that joins the two tight points,
\[
  H-48=(z-4)(12-z),\qquad g=4z(12-z),\qquad\text{so}\qquad H-48=\frac{z-4}{4z}\,g .
\]
This suggests the variable multiplier $\lambda(z)=\frac{z-4}{4z}$, which is $0$ at $z=4$ and $\frac16$ at
$z=12$. A two-line computation (checked symbolically in the script) gives the identity
\begin{equation}\label{COMP:eq:cert}
  H-48=\frac{P}{z}+\frac{z-4}{4z}\,g,\qquad
  P=\frac{xy-16}4\,z^2+\big(16-(x-8)(y-8)\big)z-4(x-y)^2 .
\end{equation}
We write $P=Az^2+Bz+C$ with $A=\frac{xy-16}4$, $B=16-(x-8)(y-8)$, $C=-4(x-y)^2$. Since $\lambda(z)\ge0$ and
$g\ge0$ on $D$, it suffices to prove
\begin{equation}\label{COMP:eq:Pclaim}
  P(x,y,z)\ge0\qquad\text{on the polyhedron } \{4\le x\le z,\ 4\le y\le z\le16\}.
\end{equation}
The curved constraint has disappeared. The price: $P$ vanishes on the whole edge $x=y=4$.

\paragraph{The search.} We prove \eqref{COMP:eq:Pclaim} by branch and bound over boxes, starting from
$[4,16]^3$ and bisecting the longest side. A box is \emph{settled} if one of three tests succeeds:
\begin{enumerate}[label=(\roman*),nosep]
\item it misses the region ($z<x$ or $z<y$ everywhere on the box);
\item \emph{value test:} the interval evaluation of $P$ on the box has lower end $\ge0$;
\item \emph{derivative test:} interval enclosures of $\partial P/\partial x$ and $\partial P/\partial y$
  (from dual numbers) have lower ends $\ge0$, so $P$ is smallest on the face $x=\underline x$, $y=\underline y$,
  and the interval evaluation of $P$ on that face has lower end $\ge0$.
\end{enumerate}
Test (iii) is what handles the tight edge: on it, $A=B=C=0$ exactly, so the evaluation on the face gives
the interval $[0,0]$. All arithmetic is in exact rationals (Python's
\texttt{fractions.Fraction}), so there is no rounding at all. Here is the heart of
\texttt{code/thue\_bb.py}, slightly abridged:
{\footnotesize
\begin{verbatim}
from fractions import Fraction as Fr
class I:                                   # closed interval, rational endpoints
    def __init__(s, lo, hi=None): s.lo, s.hi = Fr(lo), Fr(lo if hi is None else hi)
    def __add__(s, t): t = lift(t); return I(s.lo + t.lo, s.hi + t.hi)
    def __neg__(s): return I(-s.hi, -s.lo)
    def __mul__(s, t):
        t = lift(t); p = [s.lo*t.lo, s.lo*t.hi, s.hi*t.lo, s.hi*t.hi]
        return I(min(p), max(p))
    def sq(s):                             # a square is >= 0
        lo, hi = sorted([abs(s.lo), abs(s.hi)])
        return I(0 if s.lo <= 0 <= s.hi else lo*lo, hi*hi)
class D:                                   # dual number: value and gradient, all intervals
    def __init__(s, v, g): s.v, s.g = v, g
    def __add__(s, t): t = liftD(t); return D(s.v + t.v, [a + b for a, b in zip(s.g, t.g)])
    def __mul__(s, t):                     # Leibniz rule
        t = liftD(t); return D(s.v*t.v, [s.v*b + t.v*a for a, b in zip(s.g, t.g)])
    def sq(s): return s*s
def P(x, y, z):
    A = Fr(1, 4)*(x*y - 16); B = 16 - (x - 8)*(y - 8); C = -4*(x - y).sq()
    return A*z.sq() + B*z + C
def grad(X, Y, Z):
    e = lambda k: [I(1) if j == k else I(0) for j in range(3)]
    return P(D(X, e(0)), D(Y, e(1)), D(Z, e(2))).g
def check(X, Y, Z):
    if Z.hi < X.lo or Z.hi < Y.lo: return True             # (i)   outside the region
    if P(X, Y, Z).lo >= 0: return True                      # (ii)  value test
    gx, gy, _ = grad(X, Y, Z)                               # (iii) derivative test
    return gx.lo >= 0 and gy.lo >= 0 and P(I(X.lo), I(Y.lo), Z).lo >= 0
def bb(box):                    # returns only if P >= 0 is proved on the box
    if check(*box): return
    k = max(range(3), key=lambda j: box[j].hi - box[j].lo)
    m = (box[k].lo + box[k].hi) / 2
    for half in (I(box[k].lo, m), I(m, box[k].hi)):
        b = list(box); b[k] = half; bb(tuple(b))
bb((I(4, 16), I(4, 16), I(4, 16)))
\end{verbatim}
}
\noindent (Here \texttt{lift} and \texttt{liftD} turn numbers into point intervals and constant dual
numbers; subtraction and the reversed operations are defined in the same way.)

\paragraph{The result.} The program terminates in under a second. It examines $53$ boxes; the $27$ leaves
are settled by test (i) in $9$ cases, by test (ii) in $15$ and by test (iii) in $3$; the deepest box is at
depth $8$. Since every leaf passed a valid test and the leaves cover $[4,16]^3$, this proves
\eqref{COMP:eq:Pclaim}, hence $H\ge48$ on $D$, hence the inequality of
Lemma~\ref{L1:lem:triangle-density}. (The equality case is not addressed by the search.)

Without test (iii) the search never ends. With the depth capped at $3k$, the number of unsettled boxes
is $8,16,32,64,128,256$ for $k=3,\dots,8$, all of them strung along the edge $x=y=4$: each halving of the
box size doubles the work. Figure~\ref{COMP:fig:bb} shows both runs on the slice $z=6$.

\colorlet{COMPcolval}{addcol!25}\colorlet{COMPcolmono}{green!45!black!35}
\colorlet{COMPcolout}{gray!18}\colorlet{COMPcolopen}{flagcol!70}
\begin{figure}[ht]
\centering
\begin{tikzpicture}
  % panel A: certified run, slice z=6, coordinates (x,y) in [4,16]^2, 0.36 cm per unit
  \begin{scope}[x=0.36cm,y=0.36cm,shift={(-4,-4)}]
    \foreach \t/\a/\b/\c/\d in {
      mono/4/5.5/4/7, val/5.5/7/4/7, val/4/5.5/7/8.5, out/4/5.5/8.5/10, val/5.5/7/7/10, val/7/8.5/4/7,
      out/8.5/10/4/7, val/7/10/7/10, out/4/7/10/13, out/4/7/13/16, val/7/10/10/13, out/7/10/13/16,
      out/10/13/4/7, val/10/13/7/10, out/13/16/4/10, val/10/13/10/13, out/10/13/13/16, out/13/16/10/16} {
      \draw[thin,fill=COMPcol\t] (\a,\c) rectangle (\b,\d);
    }
    \draw[thick,dashed] (6,4)--(6,6)--(4,6);
    \fill[black] (4,4) circle (2.2pt);
    \draw[thick] (4,4) rectangle (16,16);
    \node[font=\scriptsize,below] at (4,4) {$4$}; \node[font=\scriptsize,below] at (16,4) {$16$};
    \node[font=\scriptsize,left] at (4,16) {$16$};
    \node[font=\small,below] at (10,3.3) {$x$}; \node[font=\small,left] at (3.2,10) {$y$};
    \node[font=\small,align=center] at (10,18) {certified run (all tests)};
  \end{scope}
  % panel B: no derivative test, depth cap 18, zoom on [4,7]^2, 1.45 cm per unit
  \begin{scope}[x=1.45cm,y=1.45cm,xshift=6.6cm,shift={(-4,-4)}]
    \clip (3.98,3.98) rectangle (7.02,7.02);
    \foreach \t/\a/\b/\c/\d in {
      open/4.0/4.1875/4.0/4.1875, val/4.0/4.1875/4.1875/4.375, val/4.1875/4.375/4.0/4.375,
      val/4.0/4.375/4.375/4.75, val/4.375/4.75/4.0/4.75, val/4.0/4.75/4.75/5.5, val/4.75/5.5/4.0/5.5,
      val/4.0/5.5/5.5/7.0, val/5.5/7.0/4.0/7.0, val/5.5/7.0/7.0/10.0, val/4.0/5.5/7.0/8.5} {
      \draw[thin,fill=COMPcol\t] (\a,\c) rectangle (\b,\d);
    }
    \draw[thick,dashed] (6,4)--(6,6)--(4,6);
    \fill[black] (4,4) circle (2.2pt);
  \end{scope}
  \begin{scope}[x=1.45cm,y=1.45cm,xshift=6.6cm,shift={(-4,-4)}]
    \draw[thick] (4,4) rectangle (7,7);
    \node[font=\scriptsize,below] at (4,4) {$4$}; \node[font=\scriptsize,below] at (7,4) {$7$};
    \node[font=\scriptsize,left] at (4,7) {$7$};
    \node[font=\small,below] at (5.5,3.83) {$x$}; \node[font=\small,left] at (3.8,5.5) {$y$};
    \node[font=\small,align=center] at (5.5,7.45) {without (iii), depth $18$ (zoom)};
  \end{scope}
\end{tikzpicture}
\\[4pt]
{\small
\tikz\draw[fill=COMPcolval] (0,0) rectangle (0.3,0.2); value test \quad
\tikz\draw[fill=COMPcolmono] (0,0) rectangle (0.3,0.2); derivative test \quad
\tikz\draw[fill=COMPcolout] (0,0) rectangle (0.3,0.2); outside the region \quad
\tikz\draw[fill=COMPcolopen] (0,0) rectangle (0.3,0.2); unsettled}
\caption{Branch and bound for \eqref{COMP:eq:Pclaim}, shown on the slice $z=6$ (the boxes are
three-dimensional; we draw their intersections with this plane). The dashed square is the region
$x,y\le z=6$, and the dot is the tight edge $x=y=4$. Left: the certified run, with $18$ boxes on this
slice. Right: without the derivative test the boxes pile up geometrically towards the tight edge, and the
corner box is never settled.}
\label{COMP:fig:bb}
\end{figure}

\begin{remark}[What was learned]
(1) The tight points dictated the method: a strict inequality would have been proved by test (ii)
alone, and the multiplier would not have been needed. Every sharp inequality in Hales's proof needed a
local argument near its equality case of exactly this kind. (2) The multiplier $\lambda(z)$ is a
Positivstellensatz-type certificate in the sense of Section~\ref{COMP:sec:verify}: \eqref{COMP:eq:cert} reads
$z(H-48)=P+\frac{z-4}4\,g$, and only the nonnegativity of $P$ on a polyhedron is left to the computer. A
semidefinite program could look for a sum-of-squares certificate for $P$ too (Exercise~\ref{COMP:exc:sos}).
(3) The derivative test is the monotonicity argument of Lecture~1 in disguise: there the area
$\frac12ac\sin B$ is increased by lengthening $a$ and $c$.
\end{remark}

%======================================================================
\section{Gallery of theorems}
\label{COMP:sec:gallery}

The theorems below were proved with essential help from computers, or are certified by finite objects
that were found by computer. Each comes with a sentence on how it fits the course.

\begin{theorem}[Hales, Ann.\ of Math.\ 162 (2005); Hales--Ferguson, Discrete Comput.\ Geom.\ 36 (2006)]
No packing of congruent balls in $\R^3$ has density greater than $\pi/\sqrt{18}$.
\end{theorem}
\noindent The formal proof is T.~Hales et al., \emph{A formal proof of the Kepler conjecture}, Forum Math.\ Pi 5
(2017), e2. Cases are indexed by tame plane graphs (Section~\ref{COMP:sec:index}) and settled by linear programs
and interval arithmetic; the whole of this interlude is a miniature of it.

\begin{theorem}[Hales--McLaughlin, J.\ Amer.\ Math.\ Soc.\ 23 (2010)]
Every Voronoi cell of a packing of unit balls in $\R^3$ has volume at least that of the regular dodecahedron
of inradius $1$.
\end{theorem}
\noindent Fejes T\'oth's dodecahedral conjecture (Interlude~E, Section~\ref{SM:sec:frustration}), proved by the
same computer-assisted methodology as the Kepler conjecture.

\begin{theorem}[Musin--Tarasov, Discrete Comput.\ Geom.\ 48 (2012); Experiment.\ Math.\ 24 (2015)]
The optimal Tammes configurations for $N=13$ and $N=14$ are unique up to isometry, with minimal angular
distances $\approx57.1367^\circ$ and $\approx55.6706^\circ$.
\end{theorem}
\noindent The proofs enumerate irreducible contact graphs: the full-scale version of the five-point example of
Section~\ref{COMP:sec:index}. The case $N=13$ is the ``strong thirteen spheres problem'': how large can $13$ equal balls touching a unit ball
be? Since $57.1367^\circ<60^\circ$, it contains a new proof of $\kiss_3=12$.

\begin{theorem}[Odlyzko--Sloane, J.\ Combin.\ Theory Ser.~A 26 (1979); Levenshtein, Soviet Math.\ Dokl.\ 20 (1979)]
$\kiss_8=240$ and $\kiss_{24}=196\,560$.
\end{theorem}
\noindent Example~\ref{COMP:ex:240}. The computer found the polynomials; the proof is checkable by hand.
Bannai and Sloane (Canad.\ J.\ Math.\ 33, 1981) proved uniqueness of the configurations.

\begin{theorem}[Musin, Ann.\ of Math.\ 168 (2008)]
$\kiss_4=24$.
\end{theorem}
\noindent The plain LP bound gives only $25$. Musin modifies Delsarte's method so that the auxiliary polynomial
may be positive on part of the range, and controls that part by a separate geometric argument.

\begin{theorem}[Bachoc--Vallentin, J.\ Amer.\ Math.\ Soc.\ 21 (2008)]
$\kiss_5\le45$, $\kiss_6\le78$, $\kiss_7\le135$, $\kiss_9\le366$, $\kiss_{10}\le567$.
\end{theorem}
\noindent Three-point semidefinite bounds, a matrix version of Theorem~\ref{COMP:thm:LP}. N.~Leijenhorst and
D.~de Laat (Math.\ Program.\ Comput., 2024) improved the bounds in dimensions $11$--$23$, and
H.~Cohn, D.~de Laat and A.~Salmon (\emph{Three-point bounds for sphere packing}, arXiv:2206.15373) gave the
analogous improvement of the Cohn--Elkies bound for packings in dimensions $4$--$7$ and $9$--$16$.

\begin{theorem}[Mark\'ot--Csendes, SIAM J.\ Optim.\ 16 (2005); Mark\'ot, J.\ Global Optim.\ 81 (2021)]
The best known packings of $n$ equal circles in a square are optimal for $n=28,\dots,33$.
\end{theorem}
\noindent Interval branch and bound in configuration space, the method of Section~\ref{COMP:sec:bb} with a
great deal of engineering.

\begin{theorem}[Brakensiek--Heule--Mackey--Narv\'aez, IJCAR 2020; J.\ Automat.\ Reason.\ 66 (2022)]
Every tiling of $\R^7$ by unit cubes contains two cubes sharing a full face (Keller's conjecture holds in
dimension $7$).
\end{theorem}
\noindent Keller's conjecture was known for $n\le6$ (Perron) and false for $n\ge8$ (Lagarias--Shor for $n\ge10$, Mackey
for $n=8$). The case $n=7$ was reduced to the non-existence of a clique of size $128$ in certain graphs and settled by
a SAT solver; the proofs of unsatisfiability (about $200$ gigabytes) were checked by an independent, formally
verified checker. A packing problem for cubes, settled by a certificate of non-existence.

\begin{theorem}[Appel--Haken 1977; Robertson--Sanders--Seymour--Thomas, J.\ Combin.\ Theory Ser.~B 70 (1997); Gonthier, Notices AMS 55 (2008)]
Every planar graph is $4$-colourable.
\end{theorem}
\noindent The model case of ``cases indexed by planar graphs'': an unavoidable set of configurations ($633$ in the
1997 proof), each shown reducible by computer. Gonthier's Coq proof removed the doubts about the code.

\begin{theorem}[Tucker, C.~R.\ Acad.\ Sci.\ Paris 328 (1999); Found.\ Comput.\ Math.\ 2 (2002)]
The Lorenz equations, with the classical parameters, have a robust strange attractor.
\end{theorem}
\noindent Smale's 14th problem, solved by interval arithmetic: an example far from geometry of the methods of
Section~\ref{COMP:sec:verify}.

\begin{theorem}[Lower bounds by construction, 2025]
$\kiss_{11}\ge593$ (AlphaEvolve, Google DeepMind, arXiv:2506.13131); $\kiss_{10}\ge510$, $\kiss_{14}\ge1932$
(M.~Ganzhinov, \emph{Highly symmetric lines}, Linear Algebra Appl.\ 722 (2025); arXiv:2207.08266).
\end{theorem}
\noindent A lower bound is certified by a list of vectors and an exact check of all pairwise inner products.
Finding the list is the hard part; checking it is easy. This is the opposite situation to the upper bounds above.

%======================================================================
\section*{Exercises}
\addcontentsline{toc}{section}{Exercises}

\begin{exercise}
Using Lemma~\ref{COMP:lem:angle}, show that in a configuration with $\rho\ge\pi/3$ every point has at most five
contacts, and that the bound is still five for the Tammes configuration of $13$ points
($\theta\approx57.14^\circ$). Show that the bound drops to four exactly when $\theta>\arctan2$, the edge of the
regular icosahedron.
\end{exercise}

\begin{exercise}
Modify the program \texttt{code/contact5.py} to run the rules R1--R4 on the $156$ graphs on six vertices.
Which graphs survive?
Find the critical configurations whose stressed graphs are the octahedron, the hexagon $C_6$, and $C_5$ plus an
isolated point. Are there others?
\end{exercise}

\begin{exercise}
Prove the equality case of Proposition~\ref{COMP:prop:rankin}(b): if $2d$ nonzero vectors in $\R^d$ have pairwise
non-positive inner products, they are $\pm$ an orthogonal basis up to positive scaling. Deduce that the octahedron
is the unique optimal configuration for $N=6$, while for $N=5$ the optimum is not unique.
\end{exercise}

\begin{exercise}
Verify Example~\ref{COMP:ex:240}: expand $f$ in Gegenbauer polynomials for $n=8$ and compute $f(1)/c_0$. Then
use the same polynomial in dimension $3$. Show that its coefficients are still non-negative, and that it only
proves $\kiss_3\le1890/127\approx14.9$. Why should one expect a polynomial with roots at $-\frac12,0,\frac12$ to
be good in dimension $8$ but not in dimension $3$?
\end{exercise}

\begin{exercise}\label{COMP:exc:kkt}
For $N=4$, compute the KKT multipliers of Proposition~\ref{COMP:prop:kkt} at the regular tetrahedron and at the
square on a great circle. Show that the square is not first-order jammed: find a nontrivial infinitesimal motion $v$ with $Av=0$, and
check that along the corresponding curve all four contact distances increase to second order. Relate this to
Example~\ref{TOP:ex:N4pass} of Interlude~D.
\end{exercise}

\begin{exercise}
(a) Compute the natural interval extension of $f(x)=x^2-2x$ on $[0,2]$, and of the equivalent forms $x(x-2)$ and
$(x-1)^2-1$. (b) Use Proposition~\ref{COMP:prop:newton} with exact rational arithmetic to enclose $\sqrt2$ in an
interval of width less than $10^{-20}$, starting from $[1,2]$.
\end{exercise}

\begin{exercise}\label{COMP:exc:sos}
(a) Check that the three coefficients $A,B,C$ of $P$ vanish on the edge $x=y=4$.
Compute $\partial P/\partial x$ and $\partial P/\partial y$, and show by interval arithmetic that both are
$\ge0$ on the box $[4,5.5]\times[4,7]\times[4,7]$, one of the three boxes settled by test (iii). (b) Using an SDP solver,
look for a certificate $P=\sigma_0+\sum_i\sigma_ig_i$ on the polyhedron of \eqref{COMP:eq:Pclaim} with
$g\in\{x-4,\,y-4,\,z-x,\,z-y,\,16-z\}$, and round it to exact rationals.
\end{exercise}

\begin{exercise}
Write the statement ``there are $5$ points on $\Sph^2$ with pairwise distances greater than $\pi/2$'' as a
first-order sentence. How many variables and polynomials does it have? Using
Proposition~\ref{COMP:prop:rankin}, decide it by hand. Why does the existence of a decision procedure imply that
$\cos d_N$ is algebraic for every $N$?
\end{exercise}

\begin{exercise}
Modify \texttt{code/e8\_mitm.py} to count the vectors of norm $2$ and $4$ in $D_4$ (split $2+2$) and in $E_7$
(take $E_7$ to be the vectors of $E_8$ orthogonal to the root $e_1-e_2$, and keep the first two coordinates in the
same half). You should find $24,24$ and $126,756$.
\end{exercise}

%======================================================================
\section*{Further reading}
\addcontentsline{toc}{section}{Further reading}
\begin{itemize}[nosep]
\item R.~E.~Moore, R.~B.~Kearfott, M.~J.~Cloud, \emph{Introduction to Interval Analysis}, SIAM, 2009. Interval
  arithmetic, interval Newton and Krawczyk, with code.
\item W.~Tucker, \emph{Validated Numerics: A Short Introduction to Rigorous Computations}, Princeton University
  Press, 2011. Includes automatic differentiation and the Lorenz attractor.
\item G.~Blekherman, P.~A.~Parrilo, R.~R.~Thomas (eds.), \emph{Semidefinite Optimization and Convex Algebraic
  Geometry}, MOS--SIAM Series on Optimization, 2013. Sums of squares, Positivstellens\"atze, SDP.
\item S.~Basu, R.~Pollack, M.-F.~Roy, \emph{Algorithms in Real Algebraic Geometry}, 2nd ed., Springer, 2006.
  Sturm sequences, cylindrical algebraic decomposition, complexity.
\item D.~Cox, J.~Little, D.~O'Shea, \emph{Ideals, Varieties, and Algorithms}, Springer UTM. Gr\"obner bases.
\item J.~H.~Conway, N.~J.~A.~Sloane, \emph{Sphere Packings, Lattices and Groups}, Springer; Ch.~9 (bounds for
  codes and sphere packings), Ch.~13 (Odlyzko--Sloane, bounds on kissing numbers) and Ch.~14 (Bannai--Sloane,
  uniqueness).
\item T.~C.~Hales et al., \emph{A formal proof of the Kepler conjecture}, Forum Math.\ Pi 5 (2017), e2; and
  T.~C.~Hales, \emph{Dense Sphere Packings: A Blueprint for Formal Proofs}, LMS Lecture Note Series 400, Cambridge,
  2012.
\item J.~C.~Lagarias (ed.), \emph{The Kepler Conjecture: The Hales--Ferguson Proof}, Springer, 2011.
\item H.~Cohn, \emph{Packing, coding, and ground states}, PCMI lecture notes, arXiv:1603.05202. LP and SDP bounds
  from the point of view of these notes.
\end{itemize}
